\chapter{Lead--lag inter-actifs}\label{ch:s1:cross-asset-lead-lag}

La même nouvelle atteint les prix de différents actifs à des vitesses différentes. Hong, Torous et Valkanov ont montré que les rendements de secteurs comme le commerce de détail,
les métaux et le pétrole prévoient l'ensemble du marché boursier américain jusqu'à deux mois à l'avance~; Hou a montré que les petites entreprises suivent les grandes entreprises de leur propre secteur, le plus
lentement après une mauvaise nouvelle~; à l'autre bout de l'échelle des temps, les contrats à terme et les actions ou fonds qui les répliquent sont liés par des avances de quelques secondes. Sur les marchés synthétiques de ce
chapitre, une avance de cinq jours des matières premières vers trois secteurs est détectée sur deux de ses trois liens par un balayage rigoureux et négociée à un ratio de Sharpe de 1.73 après coûts hors échantillon~; une avance de deux
secondes entre deux actifs est détectée à chaque séance et rapporte à un trader qui est au plus une seconde
en retard. La construction s'appuie sur \texttt{firm.xasset}.

\section{Signaux intermarchés à horizon quotidien}

\begin{definition}[Signal intermarchés, diffusion lente de l'information]\label{def:s1:cross-asset-lead-lag:intermarket}
Un \emph{signal intermarchés}\index{signal intermarchés} prédit les rendements d'un marché à partir des rendements passés d'un autre marché qui reflète plus tôt les mêmes fondamentaux,
comme une matière première pour ses producteurs ou un marché obligataire pour les actions sensibles aux taux. La \emph{diffusion lente de l'information}\index{diffusion lente de l'information} est
le mécanisme qui le sous-tend~: l'information atteint les prix de différents actifs à des vitesses différentes parce que les investisseurs se spécialisent, suivent peu de marchés ou
font face à des coûts de transaction, de sorte que des prix qui devraient évoluer ensemble évoluent successivement.
\end{definition}

Une relation d'avance--retard (lead--lag) (livre~7, chapitre~10) se détecte en corrélant une série avec le passé de l'autre. À horizon quotidien, les candidats sont nombreux et
les liens réels peu nombreux. Le chapitre les introduit volontairement~: les rendements des dix facteurs sectoriels de \texttt{firm.synthmkt} suivent les dix matières premières de \texttt{firm.synthfut}
à travers trois liens, chaque secteur recevant 0.08 fois la moyenne des rendements d'une matière première sur les cinq jours précédents
(\cref{lst:s1:cross-asset-lead-lag:scan}). L'effet est faible~: avec une volatilité des matières premières d'environ 25~\% et une volatilité sectorielle de 8~\%, le retard de chaque jour
porte une corrélation d'environ 0.05, soit une statistique de Student inférieure à 2 sur cinq ans.

Deux balayages sont menés sur les cinq premières années. Le balayage retard par retard corrèle chaque secteur avec chaque matière première à des retards de 1 à 10 jours, soit 1\,000 tests, avec un seuil
de Bonferroni de 4.06 pour une erreur familiale de 5~\%. Il détecte un seul lien, le premier introduit, au bon retard de cinq jours ($t = 4.7$)~; les deux autres liens introduits
culminent à 2.7 et 3.5, sous le seuil (\cref{fig:s1:cross-asset-lead-lag:scan}). Le balayage par fenêtre effectue une seule statistique par
couple, la corrélation de chaque secteur avec le rendement passé de la matière première sur cinq jours, soit 100 tests avec un seuil de 3.48. Il détecte les deux autres ($t = 4.1$
chacun) et manque le premier ($t = 3.2$). Aucun des deux balayages ne trouve tout, et aucun ne trouve de faux positif. Un balayage adapté au mécanisme, ici une fenêtre plutôt que des retards isolés, a
davantage de puissance, ce qui impose de poser le mécanisme avant de lancer la recherche.

\begin{omfigure}
\begin{tikzpicture}
\begin{axis}[width=10.5cm,height=5.2cm,xlabel={\small retard (jours)},ylabel={\small statistique de Student, cinq premières années},tick label style={font=\footnotesize},
  xmin=0.5,xmax=10.5,ymin=-5,ymax=5.5,grid=major,grid style={gray!25},legend columns=4,legend style={font=\scriptsize,at={(0.5,-0.3)},anchor=north,
  draw=none,fill=none,/tikz/every even column/.append style={column sep=6pt}},table/col sep=comma]
\addplot[omThm,thick,mark=*] table[x=lag,y=link1]{figdata/strategies-1/22-cross-asset-lead-lag/lagscan.csv};
\addplot[omDef,thick,mark=square*] table[x=lag,y=link2]{figdata/strategies-1/22-cross-asset-lead-lag/lagscan.csv};
\addplot[omExo,thick,mark=triangle*] table[x=lag,y=link3]{figdata/strategies-1/22-cross-asset-lead-lag/lagscan.csv};
\addplot[black!60,dashed] coordinates {(0.5,4.06) (10.5,4.06)};
\addplot[black!60,dashed,forget plot] coordinates {(0.5,-4.06) (10.5,-4.06)};
\legend{lien 1,lien 2,lien 3,Bonferroni (1\,000 tests)}
\end{axis}
\end{tikzpicture}

\omcaption{Le balayage retard par retard sur les secteurs et matières premières synthétiques, cinq premières années~: la statistique de Student de la corrélation décalée de chaque lien introduit,
pour des retards de 1 à 10 jours, comparée au seuil de Bonferroni pour 1\,000 tests. Chaque lien répartit le rendement d'une matière première sur les cinq jours suivants.
Données~: \texttt{s1\_xasset.daily}.}\label{fig:s1:cross-asset-lead-lag:scan}
\end{omfigure}

Les liens retenus sont négociés sur les cinq dernières années~: chaque secteur lié est détenu en proportion du rendement de sa matière première sur cinq jours, en écarts-types,
plafonné à deux, et rééquilibré chaque jour à 2 points de base par unité échangée. Le ratio de Sharpe hors échantillon est de 2.09 avant coûts et de 1.73 après, sur deux liens. Ce
niveau est élevé parce que la diffusion introduite est propre et constante~; le registre public est mensuel, plus faible, et varie dans le temps.

\section{Lead--lag intrajournalier}

À horizon intrajournalier, le meneur est généralement un contrat à terme et les suiveurs sont les instruments au comptant qui le répliquent (livre~7, chapitre~10,
avance du contrat à terme sur le comptant). Le chapitre simule vingt séances de 23\,400 secondes où un suiveur voit le prix efficient du meneur avec deux secondes de retard,
les deux étant observés à des instants de cotation aléatoires (une par seconde en moyenne) et avec du bruit (\cref{lst:s1:cross-asset-lead-lag:trades}). La corrélation croisée de Hayashi--Yoshida
(livre~4, chapitre~21) traite les cotations asynchrones sans grille commune~; son pic se situe à deux secondes dans chacune des vingt
séances, et sa forme est symétrique autour de l'avance (\cref{fig:s1:cross-asset-lead-lag:ccf}).

\begin{omfigure}
\begin{tikzpicture}
\begin{axis}[width=10cm,height=4.8cm,ybar,bar width=8pt,xlabel={\small retard du suiveur (secondes)},ylabel={\small corrélation croisée (pic = 1)},
  tick label style={font=\footnotesize},xmin=-6.8,xmax=6.8,ymin=0,ymax=1.1,grid=major,grid style={gray!25},table/col sep=comma]
\addplot[fill=omThm!60,draw=omThm] table[x=lag,y=ccf]{figdata/strategies-1/22-cross-asset-lead-lag/ccf.csv};
\end{axis}
\end{tikzpicture}

\omcaption{La corrélation croisée de Hayashi--Yoshida entre un meneur et un suiveur qui voit le prix du meneur avec deux secondes de retard, moyennée sur vingt
séances simulées et ramenée à un pic égal à un~; un pic à un retard positif signifie que la seconde série suit. Données~:
\texttt{s1\_xasset.intraday}.}\label{fig:s1:cross-asset-lead-lag:ccf}
\end{omfigure}

La transaction~: lorsque le meneur a bougé de plus de 2 points de base au cours des deux dernières secondes, on achète ou on vend le suiveur dans le même sens après une latence,
on le conserve deux secondes, et on paie un demi-point de base dans chaque sens.

\begin{center}
\small
\begin{tabular}{@{}lcccc@{}}
\toprule
latence (secondes) & 0 & 1 & 2 & 3\\
\midrule
transactions par séance & 2\,343 & 2\,193 & 2\,025 & 1\,862\\
rendement net par transaction (pb) & 0.73 & 0.35 & $-0.51$ & $-0.81$\\
part des transactions gagnantes & 69.6~\% & 57.2~\% & 32.0~\% & 25.0~\%\\
séances à total positif & 20 & 20 & 0 & 0\\
\bottomrule
\end{tabular}
\end{center}

Un trader qui agit immédiatement conserve l'essentiel du mouvement de deux secondes~; avec une seconde de retard, il n'en conserve que la moitié~; à deux secondes, le suiveur a déjà bougé et la
transaction paie la fourchette pour rien. L'avance est la même à toute latence. Ce qui change, c'est qui peut en profiter, comme pour les nouvelles du chapitre~17.

\section{Diffusion lente et lenteur des acteurs}

Le registre public indique d'où vient la lenteur. Menzly et Ozbas ont montré que les actions de secteurs liés par des relations fournisseur--client se prédisent
mutuellement leurs rendements, et que cette prévisibilité croisée est plus faible lorsque davantage d'investisseurs informés suivent les actions (plus d'analystes, plus de
détention institutionnelle)~: la spécialisation rend les investisseurs lents à percevoir les nouvelles des secteurs qu'ils ne couvrent pas. Hou a montré que l'avance des grandes entreprises sur les petites
est surtout intrasectorielle, provient d'un ajustement lent aux mauvaises nouvelles, et est la plus forte dans les secteurs petits, moins concurrentiels et négligés. Hong,
Torous et Valkanov~: leurs secteurs devançaient le marché de deux mois au plus, et ceux qui le faisaient étaient aussi ceux qui prévoyaient
l'activité économique. Leur note de 2014, qui étend l'échantillon jusqu'en 2013 et publie le code, trouve un noyau robuste de secteurs précurseurs mais un ensemble plus restreint que par le passé,
les relations prédictives évoluant dans le temps. Les avances s'estompent à mesure qu'elles deviennent connues, comme la plupart des signaux publiés (livre~7, chapitre~13).

\section{Mettre en œuvre}

Une stratégie intermarchés se construit en quatre étapes~: un mécanisme qui dit qui doit mener qui et à quelle vitesse~; un balayage adapté à celui-ci, avec la
correction des tests multiples pour le nombre de candidats~; un test de stabilité dans le temps, puisque les avances apparaissent et disparaissent~; et les coûts, car le
mouvement prévu est faible. Les versions quotidienne et intrajournalière diffèrent par ce qui les limite. À horizon quotidien, la limite est statistique~: quelques liens réels parmi
des centaines de candidats, chacun faible. À horizon intrajournalier, la limite est la vitesse et le coût~: l'avance est évidente, et elle revient à celui qui la voit en premier
et qui paie le moins cher pour la négocier.

\section{Fiches de stratégie}

\begin{strategyfile}{Signal obligations-actions}\label{strat:s1:cross-asset-lead-lag:bond}
\sfield{À qui profite le signal, et pourquoi} Les investisseurs en actions réagissent avec retard aux nouvelles sur les taux dans les secteurs sensibles aux taux.
\sfield{Instruments et plateformes} Contrats à terme sur obligations comme variables avancées~; banques, services publics, immobilier et autres actions ou contrats à terme sectoriels sensibles aux taux comme variables suiveuses.
\sfield{Signal} Rendements récents des contrats à terme sur obligations, projetés sur les secteurs selon leur sensibilité estimée aux taux.
\sfield{Dimensionnement et exécution} Positions sectorielles dimensionnées selon le signal et la volatilité~; quotidien.
\sfield{Coûts} Modérés~; paniers sectoriels ou contrats à terme.
\sfield{Comment elle s'éteint} Des traders intermarchés plus rapides~; sensibilités aux taux qui changent.
\sfield{Horizon, capacité, infrastructure} Quelques jours.
\sfield{Backtest honnête} Sensibilités estimées sur les seules données passées~; un balayage corrigé pour chaque secteur et chaque retard testé.
\sfield{Sources} Aucun chiffre de performance vérifié~; le mécanisme est la \omterm{def:s1:cross-asset-lead-lag:intermarket}{diffusion lente de l'information} (Hong, Torous et Valkanov, 2007).
\end{strategyfile}

\begin{strategyfile}{Signal matières premières-producteurs}\label{strat:s1:cross-asset-lead-lag:commodity}
\sfield{À qui profite le signal, et pourquoi} Les investisseurs en actions, spécialisés, suivent moins étroitement les marchés des matières premières que les traders de matières premières.
\sfield{Instruments et plateformes} Contrats à terme sur matières premières comme variables avancées~; producteurs et gros consommateurs comme variables suiveuses.
\sfield{Signal} Le rendement de la matière première sur une fenêtre calée sur la vitesse de diffusion.
\sfield{Dimensionnement et exécution} Achats sur les producteurs après des hausses des matières premières, ventes après des baisses, couverts contre le marché~; quotidien ou hebdomadaire.
\sfield{Coûts} Modérés.
\sfield{Comment elle s'éteint} Des investisseurs mieux informés couvrant les deux marchés, comme l'ont montré Menzly et Ozbas.
\sfield{Horizon, capacité, infrastructure} De quelques jours à quelques semaines.
\sfield{Backtest honnête} Les liens choisis à partir d'un mécanisme, puis testés~; périodes hors échantillon (out-of-sample).
\sfield{Sources} Hong, Torous et Valkanov (2007)~; Menzly et Ozbas (2010)~; ce chapitre~: 1.73 après coûts sur deux liens implantés.
\end{strategyfile}

\begin{strategyfile}{Avance intrajournalière du contrat à terme sur le comptant}\label{strat:s1:cross-asset-lead-lag:futcash}
\sfield{À qui profite le signal, et pourquoi} Les fournisseurs de liquidité sur les instruments au comptant dont les cotations se mettent à jour après celles des contrats à terme.
\sfield{Instruments et plateformes} Contrats à terme sur indices et ETF ou actions qui les répliquent.
\sfield{Signal} Le mouvement du contrat à terme au cours des dernières secondes.
\sfield{Dimensionnement et exécution} Ordres agressifs sur la variable suiveuse, dans le sillage de la variable avancée~; petites tailles, en grand nombre.
\sfield{Coûts} Fourchettes et frais sur chaque transaction~; l'infrastructure pour la rapidité.
\sfield{Comment elle s'éteint} Concurrence sur la latence~: une seconde de retard conserve la moitié du gain, deux secondes de retard le font disparaître.
\sfield{Horizon, capacité, infrastructure} Secondes~; colocation et flux directs.
\sfield{Backtest honnête} Horodatages tels que reçus~; cotations de la variable suiveuse, et non transactions~; modèles de file d'attente et de latence.
\sfield{Sources} Aucun chiffre de performance vérifié~; simulation de ce chapitre.
\end{strategyfile}

\begin{strategyfile}{Signal devise-exportateur}\label{strat:s1:cross-asset-lead-lag:currency}
\sfield{À qui profite le signal, et pourquoi} Les investisseurs en actions réagissent avec retard aux variations de change dans les résultats des exportateurs et des importateurs.
\sfield{Instruments et plateformes} Contrats à terme ou contrats de gré à gré sur devises comme variables avancées~; exportateurs et importateurs comme variables suiveuses.
\sfield{Signal} Rendements récents des devises, multipliés par la part estimée du chiffre d'affaires étranger de chaque entreprise.
\sfield{Dimensionnement et exécution} Acheteurs des exportateurs après l'affaiblissement de la devise nationale, vendeurs après son renforcement~; couverts contre le marché.
\sfield{Coûts} Modérés.
\sfield{Comment elle s'éteint} Couvertures de change des entreprises~; analystes plus rapides.
\sfield{Horizon, capacité, infrastructure} Quelques semaines~; données de chiffre d'affaires par segment.
\sfield{Backtest honnête} Parts de chiffre d'affaires telles que publiées à l'époque.
\sfield{Sources} Aucun chiffre de performance vérifié.
\end{strategyfile}

\section{Tutoriel~: même information, vitesses différentes}\label{tut:s1:cross-asset-lead-lag:tut}

\begin{tutorial}
\textbf{Objectif.} Implanter une diffusion des matières premières vers les secteurs et la retrouver avec deux balayages~; négocier ce qui est trouvé hors échantillon (out-of-sample)~; simuler une avance de deux secondes,
la retrouver avec l'estimateur de Hayashi--Yoshida et la négocier à quatre latences. \textbf{État final~:} le tableau et les deux figures.
\begin{tutsteps}
\item \textbf{Diffusion et balayages}~: implantation, balayage retard par retard et détecteur de Bonferroni.
\omcode{code/firm/xasset/firm_xasset.py}{26}{54}{Diffusion implantée, balayage des corrélations retardées et détection.}\label{lst:s1:cross-asset-lead-lag:scan}
\item \textbf{La transaction intrajournalière}~: agir sur le mouvement de la variable avancée après une latence.
\omcode{code/firm/xasset/firm_xasset.py}{80}{93}{Négociation de la variable suiveuse après une latence.}\label{lst:s1:cross-asset-lead-lag:trades}
\item \textbf{Exécuter} \texttt{daily()} et \texttt{intraday()}, ainsi que \texttt{fig\_xasset.py}.
\end{tutsteps}
\textbf{Pour aller plus loin.} Laisser un lien s'éteindre à mi-parcours et tester la stabilité avec des balayages glissants~; corréler deux variables avancées pour qu'un faux lien apparaisse
à travers la variable corrélée~; rendre les cotations de la variable suiveuse périmées seulement une partie du temps.
\end{tutorial}

\section{Construction~: signaux intermarchés}\label{bld:s1:cross-asset-lead-lag:xasset}

\begin{build}
\bfield{Rôle} Diffusion implantée, balayages de relations d'avance--retard (lead--lag) avec contrôle des tests multiples, signaux intermarchés et simulateur intrajournalier variable avancée--variable suiveuse.
\bfield{Interface} \texttt{diffuse(follow, lead, links, beta, lag)}, \texttt{lead\_scan(lead, follow, lags)}, \texttt{detect(tstats, n\_tests, alpha)},
\texttt{signal(lead, links, window, n\_follow)}, \texttt{simulate\_pair(seconds, lead, vol, noise, rate, rng)},
\texttt{lead\_trades(g1, g2, window, latency, hold, threshold, half\_spread)}.
\bfield{Règles} Signaux fondés sur les seuls rendements passés~; balayages corrigés pour chaque test effectué~; une seule transaction intrajournalière à la fois.
\bfield{Tests d'acceptation} \texttt{code/firm/xasset/tests/}~: diffusion et signaux vérifiés à la main, un balayage qui trouve exactement un retard implanté, une transaction
vérifiée à la main.
\bfield{Pour aller plus loin} Tests de stabilité glissants~; contrôle du taux de fausses découvertes (FDR)~; exécutions intrajournalières tenant compte de la file d'attente.
\end{build}

\begin{omsources}
\item H.~Hong, W.~Torous et R.~Valkanov, ``Do industries lead stock markets?'', \emph{Journal of Financial Economics} 83(2), 2007~; et ``Note on Do
Industries Lead Stock Markets'', 2014.
\item L.~Menzly et O.~Ozbas, ``Market segmentation and cross-predictability of returns'', \emph{Journal of Finance} 65(4), 2010.
\item K.~Hou, ``Industry information diffusion and the lead-lag effect in stock returns'', \emph{Review of Financial Studies} 20(4), 2007.
\end{omsources}

\section{Exercices}

\begin{exercise}[$\star$]\label{exo:s1:cross-asset-lead-lag:1}
Quel seuil z bilatéral donne une erreur familiale de 5~\% sur 1\,000 tests avec la correction de Bonferroni~? Sur 100~?
\end{exercise}

\begin{exercise}[$\star$]\label{exo:s1:cross-asset-lead-lag:2}
Un suiveur reçoit 0.08 fois la moyenne des cinq derniers rendements journaliers d'une variable avancée. Avec des volatilités de 25~\% pour la variable avancée et 8~\% pour la variable suiveuse, quelle
est la corrélation à un retard, et sa statistique de Student sur 1\,260 jours~? Et pour la fenêtre de cinq jours sur 1\,255 jours~?
\end{exercise}

\begin{exercise}[$\star$]\label{exo:s1:cross-asset-lead-lag:3}
Pourquoi l'estimateur de Hayashi--Yoshida convient-il aux données de tick qui n'arrivent pas aux mêmes instants~?
\end{exercise}

\begin{exercise}[$\star\star$]\label{exo:s1:cross-asset-lead-lag:4}
La transaction à une latence de deux secondes perd 0.51 points de base par transaction après un coût d'un point de base. Que gagne-t-elle avant coûts, et pourquoi
reste-t-elle positive~?
\end{exercise}

\begin{exercise}[$\star\star$]\label{exo:s1:cross-asset-lead-lag:5}
Pourquoi le balayage par fenêtres trouve-t-il des liens que le balayage retard par retard manque, et inversement~?
\end{exercise}

\begin{exercise}[$\star\star$]\label{exo:s1:cross-asset-lead-lag:6}
Deux matières premières sont fortement corrélées, mais une seule précède un secteur. Que peut signaler un balayage, et comment le distinguer~?
\end{exercise}

\begin{exercise}[$\star\star\star$]\label{exo:s1:cross-asset-lead-lag:7}
\emph{Programmation.} Placer la diffusion à 0.12 au lieu de 0.08 dans \texttt{s1\_xasset.py} et relancer \texttt{daily()}. Que trouve chaque balayage, et
chaque lien est-il réel~?
\end{exercise}

\begin{exercise}[$\star\star\star$]\label{exo:s1:cross-asset-lead-lag:8}
\emph{Trouver la faille.} «~Nous avons corrélé 500 actions avec 60 contrats à terme à des retards de 1 à 20 jours, puis négocié les 50 couples les plus forts~: ratio de Sharpe du backtest de 3.~»
\end{exercise}

\section{Problème~: même information, vitesses différentes}

\begin{problem}[{Problème du week-end --- avances au jour et à la seconde}]\label{pb:s1:cross-asset-lead-lag:1}
Les marchés synthétiques du chapitre et le témoignage public.

\medskip\noindent\textbf{Partie I --- Diffusion.}
\begin{enumerate}
\item Définir un \omterm{def:s1:cross-asset-lead-lag:intermarket}{signal intermarchés} et la \omterm{def:s1:cross-asset-lead-lag:intermarket}{diffusion lente de l'information}.
\item Qu'ont établi Hong, Torous et Valkanov, et qu'a ajouté leur note de 2014~?
\item Qu'ont établi Menzly et Ozbas au sujet de la spécialisation~?
\item Qu'a établi Hou au sujet des grandes et des petites entreprises~?
\end{enumerate}

\medskip\noindent\textbf{Partie II --- Quotidien.}
\begin{enumerate}[resume]
\item Décrire la diffusion plantée et sa force par retard.
\item Décrire les deux balayages et leurs seuils.
\item Que trouve chaque balayage~?
\item Que rapportent les liens trouvés hors échantillon~?
\end{enumerate}

\medskip\noindent\textbf{Partie III --- Intrajournalier.}
\begin{enumerate}[resume]
\item Décrire la paire simulée et l'estimateur utilisé.
\item Où se situe le pic de la corrélation croisée, et à quelle fréquence~?
\item Décrire la transaction et ses coûts.
\item Donner le rendement net par transaction à chaque latence.
\end{enumerate}

\medskip\noindent\textbf{Partie IV --- Le verdict.}
\begin{enumerate}[resume]
\item Énoncer le \emph{résultat nommé}~: l'avance détectée à chaque horizon et le ratio de Sharpe de la stratégie après coûts.
\item Qu'est-ce qui limite la stratégie quotidienne, et qu'est-ce qui limite la stratégie intrajournalière~?
\item Pourquoi le mécanisme doit-il précéder le balayage~?
\item Pourquoi les avances s'estompent-elles~?
\item Comment tester la stabilité d'une avance~?
\item Quelle fiche de stratégie exige l'infrastructure la plus rapide~?
\item Quel lien le chapitre entretient-il avec le chapitre~17~?
\item En une phrase~: à qui appartient une avance~?
\end{enumerate}
\end{problem}

\section{Questions d'entretien}

\begin{interviewq}[$\star$ \iqroles{chercheur}]\label{iq:s1:cross-asset-lead-lag:1}
Comment testeriez-vous qu'un marché en précède un autre~?
\end{interviewq}

\begin{interviewq}[$\star\star$ \iqroles{chercheur}]\label{iq:s1:cross-asset-lead-lag:2}
Vous testez 10\,000 couples et en trouvez 40 significatifs à 1~\%. Que concluez-vous~?
\end{interviewq}

\begin{interviewq}[$\star\star$ \iqroles{développeur}]\label{iq:s1:cross-asset-lead-lag:3}
Comment mesureriez-vous l'avance entre un contrat à terme et un ETF à partir de données de tick aux horodatages et aux horloges différents~?
\end{interviewq}

\begin{interviewq}[$\star\star$ \iqroles{trader}]\label{iq:s1:cross-asset-lead-lag:4}
Le profit par transaction de votre stratégie d'avance du contrat à terme sur le comptant a été divisé par deux en un mois. Qu'est-il arrivé~?
\end{interviewq}

\begin{interviewq}[$\star\star$ \iqroles{risque}]\label{iq:s1:cross-asset-lead-lag:5}
Quels sont les risques d'une stratégie intermarchés qui prend des positions sur un marché en fonction d'un autre~?
\end{interviewq}

\begin{interviewq}[$\star\star\star$ \iqroles{chercheur}]\label{iq:s1:cross-asset-lead-lag:6}
Le rendement d'un suiveur s'écrit $y_t = \beta \frac{1}{L}\sum_{k=1}^{L} x_{t-k} + e_t$ avec $x$ et $e$ indépendants. Dériver la corrélation de $y_t$ avec un seul retard
$x_{t-k}$ et avec la fenêtre $\sum_{k=1}^{L} x_{t-k}$, ainsi que le rapport des deux t-statistiques.
\end{interviewq}
